\subsubsection{Reward E: Number Grounding}
\nestindent

\textbf{What it is.}
E measures the fraction of numbers in the question that also appear in the response:
\[
R_E = \frac{|\text{numbers}_\text{response} \cap \text{numbers}_\text{question}|}{|\text{numbers}_\text{question}|}
\]
Numbers are extracted via regex ($\texttt{\textbackslash b\textbackslash d+\textbackslash .?\textbackslash d*\textbackslash b}$).
A score of 1.0 means the response uses all numbers from the question in its reasoning; 0.0 means none appear.

\textbf{Why we chose it.}
To complement D in addressing the same 82.3\% reward hacking rate: E penalizes responses that ignore the question's numbers, ensuring the model uses the actual given values rather than producing plausible-looking arithmetic with unrelated numbers.
E has universal coverage---every GSM8K question contains numbers, while $\sim$3--5\% lack proper nouns---making it a stronger necessary condition for correct problem-solving.

Uesato et al.~\cite{uesato2022solving} demonstrate that verifying intermediate reasoning steps (including use of given values) improves RL training on math tasks.
However, E has lower precision than D due to frequent number collisions---many problems use common values like 2, 3, 10---so a response about a different problem may accidentally contain matching numbers.
\nestunindent
